\ifdefined\gkver\else\def\gkver{student}\fi
\documentclass[\gkver]{gaokaozhenti}
\newcommand{\ccnum}[1]{{\fontspec{Noto Sans CJK JP}#1}}
\setlist[enumerate,3]{label=(\arabic*),ref=(\arabic*)}
\begin{document}

\chapter{2018年天津理综}

%% paperType: 理综物理部分

\begin{enumerate}
\item
%% number: 1
%% paperName: 2018年天津理综第 1 题 6 分
%% typeId: 1
%% type: 单选题
%% chapter: 原子和原子核
%% point: 人工核反应
%% method: 无
%% score: 6
%% degree: 850
%% duplicateId: 0
%% body:
国家大科学过程——中国散裂中子源（CSNS）于 2017 年 8 月 28 日首次打靶成功，获得中子束流，可以为诸多领域的研究和工业应用提供先进的研究平台。下列核反应中放出的粒子为中子的是\xzanswer{B}

\onepicture[w=3.2cm]{figs/01.jpg}

\fourchoices[answer=B]
{\ce{^14_7 N} 俘获一个 $\alpha$ 粒子，产生 \ce{^17_8 O} 并放出一个粒子}
{\ce{^27_13 Al} 俘获一个 $\alpha$ 粒子，产生 \ce{^30_15 P} 并放出一个粒子}
{\ce{^11_5 B} 俘获一个质子，产生 \ce{^8_4 Be} 并放出一个粒子}
{\ce{^6_3 Li} 俘获一个质子，产生 \ce{^3_2 He} 并放出一个粒子}

%% answer: B
%% memo:
\memoanswer{
由质量数和电荷数守恒可知四个核反应方程分别为

$\ce{^14_7 N + ^4_2 He -> ^17_8 O + ^1_1 H}$

$\ce{^27_13 Al + ^4_2 He -> ^30_15 P + ^1_0 n}$

$\ce{^11_5 B + ^1_1 H -> ^8_4 Be + ^4_2 He}$

$\ce{^6_3 Li + ^1_1 H -> ^3_2 He + ^4_2 He}$

故只有 B 选项符合题意，ACD 错误。
}

\item
%% number: 2
%% paperName: 2018年天津理综第 2 题 6 分
%% typeId: 1
%% type: 单选题
%% chapter: 曲线运动
%% point: 竖直平面圆周运动
%% method: 无
%% score: 6
%% degree: 750
%% duplicateId: 0
%% body:
滑雪运动深受人民群众的喜爱，某滑雪运动员（可视为质点）由坡道进入竖直面内的圆弧形滑道 AB，从滑道的 A 点滑行到最低点 B 的过程中，由于摩擦力的存在，运动员的速率不变，则运动员沿 AB 下滑过程中\xzanswer{C}

\onepicture[w=5.0cm]{figs/02.png}

\fourchoices[answer=C]
{所受合外力始终为零}
{所受摩擦力大小不变}
{合外力做功一定为零}
{机械能始终保持不变}

%% answer: C
%% memo:
\memoanswer{
依据题意，因为运动员是做曲线运动，所以合力一定不为零，故 A 错误；对运动员进行受力分析，如图所示。

\onepicture[w=4.0cm]{figs/02_an.jpg}

由圆周运动规律可知：重力沿垂直曲面的分力与曲面对运动员的支持力的合力充当向心力，则有 $F_{N} - mg\cos\theta = m\frac{v^{2}}{R}$，即 $F_{N} = m\frac{v^{2}}{R} + mg\cos\theta$，由于运动过程中运动员速率恒定，且 $\theta$ 在减小，所以曲面对运动员的支持力越来越大，由 $f = \mu F_{N}$ 可知摩擦力越来越大，故 B 错误；由于运动员运动过程中速率不变，质量不变，动能不变，则动能变化量为零，由动能定理可得合力做功为零，故 C 正确；因为运动员运动过程中要克服摩擦力做功，所以其机械能不守恒，故 D 错误。
}

\item
%% number: 3
%% paperName: 2018年天津理综第 3 题 6 分
%% typeId: 1
%% type: 单选题
%% chapter: 电场
%% point: 带电粒子的轨迹类问题
%% method: 无
%% score: 6
%% degree: 550
%% duplicateId: 0
%% body:
如图所示，实线表示某电场的电场线（方向未标出），虚线是一带负电的粒子只在电场力作用下的运动轨迹，设 M 点和 N 点的电势分别为 $\varphi_{M}$、$\varphi_{N}$，粒子在 M 和 N 时加速度大小分别为 $a_{M}$、$a_{N}$，速度大小分别为 $v_{M}$、$v_{N}$，电势能分别为 $E_{pM}$、$E_{pN}$。下列判断正确的是\xzanswer{D}

\onepicture[w=5.0cm]{\includesvg{figs/03}}

\fourchoices[answer=D]
{$v_{M} < v_{N}$，$a_{M} < a_{N}$}
{$v_{M} < v_{N}$，$\varphi_{M} < \varphi_{N}$}
{$\varphi_{M} < \varphi_{N}$，$E_{pM} < E_{pN}$}
{$a_{M} < a_{N}$，$E_{pM} < E_{pN}$}

%% answer: D
%% memo:
\memoanswer{
依据电场线概念，可知电场线越密的地方，电场强度越大，同一个粒子在此处受到的电场力就越大，由牛顿第二定律可知其加速度就越大，则有 $a_{N} > a_{M}$；若粒子从 M 运动到 N 点，则由带电粒子所受电场力指向运动轨迹弯曲的内侧，可知在某点的电场力方向和速度方向如图所示，

\onepicture[w=4.0cm]{figs/03_an1.jpg}

故电场力做负功，电势能增大，动能减小，即 $v_{M} > v_{N}$，$E_{pM} < E_{pN}$，负电荷在低电势处电势能大，$\varphi_{M} > \varphi_{N}$；

若粒子从 N 运动到 M，则由带电粒子所受电场力指向轨迹弯曲的内侧，可知在某点的电场力方向和速度方向如图所示，

\onepicture[w=4.0cm]{figs/03_an2.jpg}

故电场力做正功，电势能减小，动能增大，即 $v_{M} > v_{N}$，$E_{pM} < E_{pN}$，负电荷在低电势处电势能大，$\varphi_{M} > \varphi_{N}$。综上所述，故 D 正确。
}

\item
%% number: 4
%% paperName: 2018年天津理综第 4 题 6 分
%% typeId: 1
%% type: 单选题
%% chapter: 交流电
%% point: 变压器
%% method: 无
%% score: 6
%% degree: 550
%% duplicateId: 0
%% body:
教学用发电机能够产生正弦式交变电流。利用该发电机（内阻可忽略）通过理想变压器向定值电阻 $R$ 供电，电路如图所示，理想交流电流表 A、理想交流电压表 V 的读数分别为 $I$、$U$，$R$ 消耗的功率为 $P$。若发电机线圈的转速变为原来的 $\frac{1}{2}$，则\xzanswer{B}

\onepicture[w=5.5cm]{figs/04.png}

\fourchoices[answer=B]
{$R$ 消耗的功率变为 $\frac{1}{2}P$}
{电压表 V 的读数为 $\frac{1}{2}U$}
{电流表 A 的读数变为 $2I$}
{通过 $R$ 的交变电流频率不变}

%% answer: B
%% memo:
\memoanswer{
由 $\omega = 2\pi n$ 可知，转速变为原来的 $\frac{1}{2}$，则角速度变为原来的 $\frac{1}{2}$；由 $E_{m} = nBS\omega$ 可知发电机产生的最大电动势为原来的 $\frac{1}{2}$，由 $U = \frac{E_{m}}{\sqrt{2}}$ 可知发电机的输出电压有效值变为原来的 $\frac{1}{2}$，即原线圈的输出电压变为原来的 $\frac{1}{2}$，由 $\frac{n_1}{n_2} = \frac{U_1}{U_2}$ 可知副线圈的输入电压变为原来的 $\frac{1}{2}$，即电压表示数变为原来的 $\frac{1}{2}$；由 $P = \frac{U^{2}}{R}$ 可知 $R$ 消耗的电功率变为 $\frac{1}{4}P$，故 A 错误 B 正确；

副线圈中的电流为 $I_2 = \frac{U}{2R}$，即变为原来的 $\frac{1}{2}$，由 $\frac{n_1}{n_2} = \frac{I_2}{I_1}$ 可知原线圈中的电流也变为原来的 $\frac{1}{2}$，故 C 错误；转速减小为原来的 $\frac{1}{2}$，则频率变为原来的 $\frac{1}{2}$，故 D 错误。
}

\item
%% number: 5
%% paperName: 2018年天津理综第 5 题 6 分
%% typeId: 1
%% type: 单选题
%% chapter: 原子和原子核
%% point: 玻尔模型
%% method: 无
%% score: 6
%% degree: 550
%% duplicateId: 0
%% body:
氢原子光谱在可见光区域内有四条谱线 $H_{\alpha}$、$H_{\beta}$、$H_{\gamma}$、$H_{\delta}$，都是氢原子中电子从量子数 $n > 2$ 的能级跃迁到 $n = 2$ 的能级发出的光，它们在真空中的波长由长到短，可以判定\xzanswer{A}

\fourchoices[answer=A]
{$H_{\alpha}$ 对应的前后能级之差最小}
{同一介质对 $H_{\alpha}$ 的折射率最大}
{同一介质中 $H_{\delta}$ 的传播速度最大}
{用 $H_{\gamma}$ 照射某一金属能发生光电效应，则 $H_{\beta}$ 也一定能}

%% answer: A
%% memo:
\memoanswer{
依据氢原子光谱规律，波长越大，频率越小，故 $H_{\alpha}$ 的频率最小，由 $E = h\nu$ 可知 $H_{\alpha}$ 对应的能量最小，由 $h\nu = E_m - E_n$ 可知 $H_{\alpha}$ 对应的前后能级之差最小，故 A 正确；$H_{\alpha}$ 的频率最小，同一介质对应的折射率最小，由 $v = \frac{c}{n}$ 可知 $H_{\alpha}$ 的传播速度最大，故 BC 错误；$H_{\gamma}$ 的波长小于 $H_{\beta}$ 的波长，故 $H_{\gamma}$ 的频率大于 $H_{\beta}$ 的频率，若用 $H_{\gamma}$ 照射某一金属能发生光电效应，则 $H_{\beta}$ 不一定能，故 D 错误。
}

\item
%% number: 6
%% paperName: 2018年天津理综第 6 题 6 分
%% typeId: 2
%% type: 多选题
%% chapter: 曲线运动
%% point: 人造地球卫星
%% method: 无
%% score: 6
%% degree: 550
%% duplicateId: 0
%% body:
2018 年 2 月 2 日，我国成功将电磁监测试验卫星“张衡一号”发射升空，标志我国成为世界上少数拥有在轨运行高精度地球物理场探测卫星的国家之一。通过观测可以得到卫星绕地球运动的周期，并已知地球的半径和地球表面的重力加速度。若将卫星绕地球的运动看作是匀速圆周运动，且不考虑地球自转的影响，根据以上数据可以计算出卫星的\xzanswer{CD}

\onepicture[w=2.4cm]{figs/06.jpg}

\fourchoices[answer=CD]
{密度}
{向心力的大小}
{离地高度}
{线速度的大小}

%% answer: CD
%% memo:
\memoanswer{
由题意，已知卫星运动的周期 $T$、地球的半径 $R$、地球表面的重力加速度 $g$，由卫星受到的万有引力提供向心力，则有 $G\frac{Mm}{r^{2}} = m\frac{4\pi^{2}}{T^{2}}r$，卫星的质量被抵消，则不能计算卫星的密度，更不能计算卫星的向心力大小，故 AB 错误；由 $G\frac{Mm}{r^{2}} = m\frac{4\pi^{2}}{T^{2}}r$，解得 $r = \sqrt[3]{\frac{GMT^{2}}{4\pi^{2}}}$，而 $r = R + h$，所以可计算卫星距离地球表面的高度，故 C 正确；由公式 $v = \frac{2\pi r}{T}$，轨道半径可以求出，周期已知，故可以计算出卫星绕地球运动的线速度，故 D 正确。
}

\item
%% number: 7
%% paperName: 2018年天津理综第 7 题 6 分
%% typeId: 2
%% type: 多选题
%% chapter: 力物体的平衡
%% point: 力的分解
%% method: 无
%% score: 6
%% degree: 550
%% duplicateId: 0
%% body:
明朝谢肇淛《五杂组》中记载：“明姑苏虎丘寺庙倾侧，议欲正之，非万缗不可。一游僧见之，曰：无烦也，我能正之。”游僧每天将木楔从塔身倾斜一侧的砖缝间敲进去，经月余扶正了塔身。假设所用的木楔为等腰三角形，木楔的顶角为 $\theta$，现在木楔背上加一力 $F$，方向如图所示，木楔两侧产生推力 $F_{N}$，则\xzanswer{BC}

\onepicture[w=5.0cm]{\includesvg{figs/07}}

\fourchoices[answer=BC]
{若 $F$ 一定，$\theta$ 大时 $F_{N}$ 大}
{若 $F$ 一定，$\theta$ 小时 $F_{N}$ 大}
{若 $\theta$ 一定，$F$ 大时 $F_{N}$ 大}
{若 $\theta$ 一定，$F$ 小时 $F_{N}$ 大}

%% answer: BC
%% memo:
\memoanswer{
由题意，选木楔为研究对象，对其进行受力分析，则木楔受到的力有：水平向左的 $F$，和两侧给它的与木楔的斜面垂直的弹力。由于木楔处于平衡状态，所以两侧给它的与木楔的斜面垂直的弹力与 $F$ 沿两侧分解的推力是相等的，力 $F$ 的分解如图所示。

\onepicture[w=4.5cm]{figs/07_an.jpg}

$F = F_1\cos\left(90^\circ - \frac{\theta}{2}\right) + F_2\cos\left(90^\circ - \frac{\theta}{2}\right) = 2F_1\cos\left(90^\circ - \frac{\theta}{2}\right) = 2F_1\sin\frac{\theta}{2}$，$F_N = F_1 = F_2$，故解得 $F_N = \frac{F}{2\sin\frac{\theta}{2}}$，所以 $F$ 一定时，$\theta$ 越小，$F_N$ 越大；$\theta$ 一定时，$F$ 越大，$F_N$ 越大，故 BC 正确 AD 错误。
}

\item
%% number: 8
%% paperName: 2018年天津理综第 8 题 6 分
%% typeId: 2
%% type: 多选题
%% chapter: 机械振动
%% point: 振幅周期频率
%% method: 无
%% score: 6
%% degree: 450
%% duplicateId: 0
%% body:
一振子沿 $x$ 轴做简谐运动，平衡位置在坐标原点。$t = 0$ 时振子的位移为 $-0.1\Um$，$t = 1\Us$ 时位移为 $0.1\Um$，则\xzanswer{AD}

\fourchoices[answer=AD]
{若振幅为 $0.1\Um$，振子的周期可能为 $\frac{2}{3}\Us$}
{若振幅为 $0.1\Um$，振子的周期可能为 $\frac{4}{5}\Us$}
{若振幅为 $0.2\Um$，振子的周期可能为 $4\Us$}
{若振幅为 $0.2\Um$，振子的周期可能为 $6\Us$}

%% answer: AD
%% memo:
\memoanswer{
若振幅为 $0.1\Um$，由题意可知从 $t = 0\Us$ 到 $t = 1\Us$ 振子经历的周期为 $\left(n + \frac{1}{2}\right)T$，则 $\left(n + \frac{1}{2}\right)T = 1\Us$（$n = 0, 1, 2, 3\cdots$），解得 $T = \frac{2}{2n+1}\Us$（$n = 0, 1, 2, 3\cdots$）。当 $n = 1$ 时，$T = \frac{2}{3}\Us$，无论 $n$ 为何值，$T$ 都不会等于 $\frac{4}{5}\Us$，故 A 正确 B 错误；

如果振幅为 $0.2\Um$，结合位移时间关系图象，有 $1\Us = \frac{T}{2} + nT\cdots$①，或者 $1\Us = \frac{5T}{6} + nT\cdots$②，或者 $1\Us = \frac{T}{6} + nT\cdots$③。对于①式，只有当 $n = 0$ 时，$T = 2\Us$，为整数；对于②式，$T$ 不为整数；对于③式，当 $n = 0$ 时，$T = 6\Us$，之后只会小于 $6\Us$，故 C 错误 D 正确。
}

\item
%% number: 9
%% paperName: 2018年天津理综第 9 题 3 分
%% typeId: 4
%% type: 实验题
%% chapter: 电路
%% point: 伏安法测电阻
%% method: 无
%% score: 3
%% degree: 650
%% duplicateId: 0
%% body:
\begin{enumerate}
	\item
	质量为 $0.45\Ukg$ 的木块静止在光滑水平面上，一质量为 $0.05\Ukg$ 的子弹以 $200\Ums$ 的水平速度击中木块，并留在其中，整个木块沿子弹原方向运动，则木块最终速度的大小是\tkanswer{20}$\Ums$。若子弹在木块中运动时受到的平均阻力为 $4.5 \times 10^{3}\UN$，则子弹射入木块的深度为\tkanswer{0.2}$\Um$。
	\item
	某研究小组做“验证力的平行四边形定则”的实验，所有器材有：方木板一块，白纸，量程为 $5\UN$ 的弹簧测力计两个，橡皮条（带两个较长的细绳套），刻度尺，图钉（若干个）。

	\begin{enumerate}
		\item
		具体操作前，同学们提出了如下关于实验操作的建议，其中正确的是\tkanswer{BC}。

		A．橡皮条应和两绳套夹角的角平分线在一条直线上

		B．重复实验再次进行验证时，结点 O 的位置可以与前一次不同

		C．使用测力计时，施力方向应沿测力计轴线；读数时视线应正对测力计刻度

		D．用两个测力计互成角度拉橡皮条时的拉力必须都小于只用一个测力计时的拉力
		\item
		该小组的同学用同一套器材做了四次实验，白纸上留下的标注信息有结点位置 O，力的标度、分力和合力的大小及表示力的作用线的点，如下图所示。其中对于提高实验精度最有利的是\tkanswer{B}。

		\fourchoices[answer=B, ispicture=true, h=3.0cm]
		{figs/09a.png}
		{figs/09b.png}
		{figs/09c.png}
		{figs/09d.png}
	\end{enumerate}
	\item
	某同学用伏安法测定待测电阻 $R_{x}$ 的阻值（约为 $10\UkO$），除了 $R_{x}$、开关 S、导线外，还有下列器材供选用：

	A．电压表（量程 $0 \sim 1\UV$，内阻约为 $10\UkO$）

	B．电压表（量程 $0 \sim 10\UV$，内阻约为 $100\UkO$）

	C．电流表（$0 \sim 1\UmA$，内阻约为 $30\UO$）

	D．电流表（$0 \sim 0.6\UA$，内阻约为 $0.05\UO$）

	E．电源（电动势 $1.5\UV$，额定电流 $0.5\UA$，内阻不计）

	F．电源（电动势 $12\UV$，额定电流 $2\UA$，内阻不计）

	G．滑动变阻器 $R_{0}$（阻值范围 $0 \sim 10\UO$，额定电流 $2\UA$）

	\begin{enumerate}
		\item
		为使测量尽量准确，电压表选用\tkanswer{B}，电流表选用\tkanswer{C}，电源选用\tkanswer{F}。（均填器材的字母代号）
		\item
		画出测量 $R_{x}$ 阻值的实验电路图\tkanswer{如图所示}。
		\item
		该同学选择器材、连接电路和操作均正确，从实验原理上看，待测电阻测量值会\tkanswer{大于}其真实值（填“大于”“小于”或“等于”），原因是\tkanswer{电压表的读数大于待测电阻两端实际电压}。
	\end{enumerate}
\end{enumerate}

%% answer: 
\jdanswer*{
\begin{enumerate}
	\item
	20，0.2

	\item
	\begin{enumerate}
		\item
		BC

		\item
		B

	\end{enumerate}

	\item
	\begin{enumerate}
		\item
		B，C，F

		\item
		见解析

		\item
		大于，电压表的读数大于待测电阻两端实际电压（其他正确表述也可）

	\end{enumerate}

\end{enumerate}
}
%% memo:
\memoanswer{
（1）由题意，设子弹的初速度为 $v_{0}$，最终速度为 $v$，子弹在木块中运动时受到的平均阻力为 $f$，子弹射入木块的深度为 $d$，由动量守恒得

$mv_{0} = (M + m)v$

代入数据，解得 $v = \frac{mv_{0}}{M + m} = \frac{0.05 \times 200}{0.45 + 0.05}\Ums = 20\Ums$，

系统减少的动能转化为克服阻力产生的内能，故有

$fd = \frac{1}{2}mv_{0}^{2} - \frac{1}{2}(M + m)v^{2}$

解得 $d = \frac{\frac{1}{2}mv_{0}^{2} - \frac{1}{2}(M + m)v^{2}}{f} = 0.2\Um$。

（2）①为了使实验结果更具有普遍性，在实验过程中不应让橡皮条的拉力方向具有特殊的角度或位置，故 A 错误；只有每一次实验时用一个力和用两个力拉的效果完全相同，即前后两次 O 点位置相同，不需要每次实验的 O 点位置都相同，故 B 正确；使用测力计时，施力方向应沿测力计轴线，读数时，视线应正对弹簧测力计的刻度，故 C 正确；合力可以大于任一个分力，也可以等于分力，还可以小于任一分力，故 D 错误。

②A、D 实验中选取的力的标度过大，导致画力时，线段太短，不利于实验精确度；B 图和 C 图选用的标度相同，但 C 中力太小，不利于作平行四边形，故 B 符合题意。

（3）①若选用电源 $1.5\UV$，由于被测电阻很大，电路中电流非常小，不利于实验测量，则电源选用 $12\UV$ 的，即 F；则电压表就应该选取 B；电路中的最大电流为 $I = \frac{12}{10000}\UA = 1.2\UmA$，故选用电流表 C。

②因为给的滑动变阻器的最大阻值为 $10\UO$，若采用限流接法，则电路中电流和电压变化不明显，故采用滑动变阻器的分压接法；由于 $\frac{R_{V}}{R_{x}} < \frac{R_{x}}{R_{A}}$，所以采用电流表内接法，则设计电路图如图所示。

\onepicture[w=4.5cm]{figs/09_an.jpg}

③由于电流表的分压作用，导致电压测量值偏大，而电流准确，由 $R_{x} = \frac{U}{I}$ 可知测量值偏大。
}

\item
%% number: 10
%% paperName: 2018年天津理综第 10 题 16 分
%% typeId: 6
%% type: 计算题
%% chapter: 运动和力
%% point: 牛顿第二定律
%% method: 无
%% score: 16
%% degree: 750
%% duplicateId: 0
%% body:
我国自行研制、具有完全自主知识产权的新一代大型喷气式客机 C919 首飞成功后，拉开了全面试验试飞的新征程。假设飞机在水平跑道上的滑跑是初速度为零的匀加速直线运动，当位移 $x = 1.6 \times 10^{3}\Um$ 时才能达到起飞所要求的速度 $v = 80\Ums$，已知飞机质量 $m = 7.0 \times 10^{4}\Ukg$，滑跑时受到的阻力为自身重力的 0.1 倍，重力加速度取 $g = 10\Umsq$，求飞机滑跑过程中
\begin{enumerate}
	\item
	加速度 $a$ 的大小；
	\item
	牵引力的平均功率 $P$。
\end{enumerate}
\onepicture[w=4.5cm, align=right]{figs/10.jpg}

%% answer: 
\jdanswer{
\begin{enumerate}
	\item
	$a = 2\Umsq$

	\item
	$P = 8.4 \times 10^{6}\UW$

\end{enumerate}
}
%% memo:
\memoanswer{
（1）飞机滑跑过程中做初速度为零的匀加速直线运动，由运动学公式得 $v^{2} = 2ax$，代入数据解得 $a = 2\Umsq$。

（2）设飞机滑跑受到的阻力为 $f$，由题意可得 $f = 0.1mg$。

设发动机的牵引力为 $F$，由牛顿第二定律得

$F - f = ma$

设飞机滑跑过程中的平均速度为 $\overline{v}$，有

$\overline{v} = \frac{v}{2}$

在滑跑阶段，牵引力的平均功率为 $P = F\overline{v}$，

$P = (ma + 0.1mg) \times \frac{v}{2}$，代入数据解得 $P = 8.4 \times 10^{6}\UW$。
}

\item
%% number: 11
%% paperName: 2018年天津理综第 11 题 18 分
%% typeId: 6
%% type: 计算题
%% chapter: 磁场
%% point: 带电粒子在磁场中的运动
%% method: 无
%% score: 18
%% degree: 450
%% duplicateId: 0
%% body:
如图所示，在水平线 ab 下方有一匀强电场，电场强度为 $E$，方向竖直向下，ab 的上方存在匀强磁场，磁感应强度为 $B$，方向垂直纸面向里，磁场中有一内、外半径分别为 $R$、$\sqrt{3}R$ 的半圆环形区域，外圆与 ab 的交点分别为 M、N。一质量为 $m$、电荷量为 $q$ 的带负电粒子在电场中 P 点静止释放，由 M 进入磁场，从 N 射出，不计粒子重力。
\begin{enumerate}
	\item
	求粒子从 P 到 M 所用的时间 $t$；
	\item
	若粒子从与 P 同一水平线上的 Q 点水平射出，同样能由 M 进入磁场，从 N 射出，粒子从 M 到 N 的过程中，始终在环形区域中运动，且所用的时间最少，求粒子在 Q 时速度 $v_{0}$ 的大小。
\end{enumerate}
\onepicture[w=6.0cm, align=right]{figs/11.png}

%% answer: 
\jdanswer{
\begin{enumerate}
	\item
	$t = \frac{\sqrt{3}RB}{E}$

	\item
	$v_{0} = \frac{qBR}{m}$

\end{enumerate}
}
%% memo:
\memoanswer{
（1）设粒子在磁场中运动的速度大小为 $v$，由洛伦兹力提供向心力得

$qvB = m\frac{v^{2}}{\sqrt{3}R}$①

设粒子在电场中运动所受电场力为 $F$，有

$F = qE$②

设粒子在电场中运动的加速度为 $a$，由牛顿第二定律得

$F = ma$③

粒子在电场中做初速度为零的匀加速直线运动，有

$v = at$④

联立①②③④式得 $t = \frac{\sqrt{3}RB}{E}$。

（2）粒子进入匀强磁场后做匀速圆周运动，其周期与速度、半径无关，运动时间只由粒子所通过的圆弧所对的圆心角的大小决定，故当轨迹与内圆相切时，所用的时间最短。设粒子在磁场中的轨迹半径为 $r'$，由几何关系可知

\onepicture[w=5.0cm]{figs/11_an.jpg}

$(r' - R)^{2} + (\sqrt{3}R)^{2} = r'^{2}$⑤

设粒子进入磁场时速度方向与 ab 的夹角为 $\theta$，即圆弧所对圆心角的一半，由几何关系可知

$\tan\theta = \frac{\sqrt{3}R}{r' - R}$⑥

粒子从 Q 射出后在电场中做类平抛运动，在电场方向上的分运动和从 P 释放后的运动情况相同，所以粒子进入磁场时沿竖直方向的速度同样为 $v$，在垂直于电场方向的分速度始终为 $v_{0}$，由运动的合成和分解可知

$\tan\theta = \frac{v}{v_{0}}$⑦

联立①⑤⑥⑦式得 $v_{0} = \frac{qBR}{m}$。
}

\item
%% number: 12
%% paperName: 2018年天津理综第 12 题 20 分
%% typeId: 6
%% type: 计算题
%% chapter: 电磁感应
%% point: 法拉第电磁感应定律
%% method: 无
%% score: 20
%% degree: 250
%% duplicateId: 0
%% body:
真空管道超高速列车的动力系统是一种将电能直接转换成平动动能的装置。图 1 是某种动力系统的简化模型，图中粗实线表示固定在水平面上间距为 $l$ 的两条平行光滑金属导轨，电阻忽略不计，ab 和 cd 是两根与导轨垂直，长度均为 $l$，电阻均为 $R$ 的金属棒，通过绝缘材料固定在列车底部，并与导轨良好接触，其间距也为 $l$，列车的总质量为 $m$。列车启动前，ab、cd 处于磁感应强度为 $B$ 的匀强磁场中，磁场方向垂直于导轨平面向下，如图 1 所示，为使列车启动，需在 M、N 间连接电动势为 $E$ 的直流电源，电源内阻及导线电阻忽略不计，列车启动后电源自动关闭。
\begin{enumerate}
	\item
	要使列车向右运行，启动时图 1 中 M、N 哪个接电源正极，并简要说明理由；
	\item
	求刚接通电源时列车加速度 $a$ 的大小；
	\item
	列车减速时，需在前方设置如图 2 所示的一系列磁感应强度为 $B$ 的匀强磁场区域，磁场宽度和相邻磁场间距均大于 $l$。若某时刻列车的速度为 $v_{0}$，此时 ab、cd 均在无磁场区域，试讨论：要使列车停下来，前方至少需要多少块这样的有界磁场？
\end{enumerate}
\twopicture[numstyle=arabic, numA=1, numB=2, align=right, wA=4.5cm, wB=7.5cm]
{figs/12a.png}
{figs/12b.png}

%% answer: 
\jdanswer{
\begin{enumerate}
	\item
	M 接电源正极，列车要向右运动，安培力方向应向右，根据左手定则，接通电源后，金属棒中电流方向由 a 到 b，由 c 到 d，故 M 接电源正极。

	\item
	$a = \frac{2BEl}{mR}$

	\item
	$\frac{I_{\text{总}}}{I_0} = \frac{mv_{0}R}{B^{2}l^{3}}$。若 $\frac{I_{\text{总}}}{I_0}$ 恰好为整数，设其为 $n$，则需设置 $n$ 块有界磁场；若 $\frac{I_{\text{总}}}{I_0}$ 不是整数，设其整数部分为 $N$，则需设置 $N + 1$ 块有界磁场。

\end{enumerate}
}
%% memo:
\memoanswer{
（1）M 接电源正极，列车要向右运动，安培力方向应向右，根据左手定则，接通电源后，金属棒中电流方向由 a 到 b，由 c 到 d，故 M 接电源正极。

（2）由题意，启动时 ab、cd 并联，设回路总电阻为 $R_{\text{总}}$，由电阻的串并联知识得

$R_{\text{总}} = \frac{R}{2}$①

设回路总电流为 $I$，由闭合电路欧姆定律得

$I = \frac{E}{R_{\text{总}}}$②

设两根金属棒所受安培力之和为 $F$，有

$F = BIl$③

由牛顿第二定律得

$F = ma$④

联立①②③④式得 $a = \frac{2BEl}{mR}$⑤

（3）设列车减速时，cd 进入磁场后经 $\Delta t$ 时间 ab 恰好进入磁场，此过程中穿过两金属棒与导轨所围回路的磁通量的变化为 $\Delta\Phi$，平均感应电动势为 $E_1$，由法拉第电磁感应定律有

$E_1 = \frac{\Delta\Phi}{\Delta t}$⑥

其中 $\Delta\Phi = Bl^{2}$⑦

设回路中平均电流为 $I'$，由闭合电路欧姆定律得

$I' = \frac{E_1}{2R}$⑧

设 cd 受到的平均安培力为 $F'$，有

$F' = BI'l$⑨

以向右为正方向，设 $\Delta t$ 时间内 cd 受安培力冲量为 $I_{\text{冲}}$，有

$I_{\text{冲}} = -F'\Delta t$⑩

同理可知，回路出磁场时 ab 受安培力冲量仍为上述值，设回路进出一块有界磁场区域安培力冲量为 $I_0$，有

$I_0 = 2I_{\text{冲}}$\ccnum{⑪}

设列车停下来受到的总冲量为 $I_{\text{总}}$，由动量定理得

$I_{\text{总}} = 0 - mv_{0}$\ccnum{⑫}

联立⑥⑦⑧⑨⑩\ccnum{⑪}\ccnum{⑫}式得 $I_{\text{总}} = mv_{0}$，$I_0 = \frac{B^{2}l^{3}}{R}$，

所以 $\frac{I_{\text{总}}}{I_0} = \frac{mv_{0}R}{B^{2}l^{3}}$。

若 $\frac{I_{\text{总}}}{I_0}$ 恰好为整数，设其为 $n$，则需设置 $n$ 块有界磁场；若 $\frac{I_{\text{总}}}{I_0}$ 不是整数，设 $\frac{I_{\text{总}}}{I_0}$ 的整数部分为 $N$，则需设置 $N + 1$ 块有界磁场。
}

\end{enumerate}

\end{document}
